% ================================================================
%  Situation Réelle 3 — Conception complète CVC d'un immeuble résidentiel R+4
%  BTS FED — Option GCF — Semestre 3 / Fin d'année
%  Chapitres mobilisés : TOUS (00 à 11) — Projet transversal
% ================================================================

\chapter{SR-3 — Conception complète CVC d'un immeuble résidentiel R+4}

\begin{tcolorbox}[
  enhanced, colback=FedGreen!8, colframe=FedGreen, arc=3pt,
  fonttitle=\bfseries\sffamily\large, coltitle=white, colbacktitle=FedGreen,
  title={Situation Réelle 3 — Fiche d'identité du projet},
  width=\textwidth
]
\begin{tabular}{@{}ll@{\hspace{2cm}}ll@{}}
  \textbf{Semestre} & 3 (fin d'année) &
  \textbf{Durée} & 8 semaines \\[4pt]
  \textbf{Groupe} & 2 à 3 étudiants &
  \textbf{Chapitres} & 00 à 11 (transversal) \\[4pt]
  \textbf{Note 1} & Dossier complet (coeff.\ 3) &
  \textbf{Note 2} & Soutenance commission (coeff.\ 2) \\[4pt]
  \textbf{Compétences} & \textbf{Toutes} C1 à C9 & & \\
\end{tabular}
\end{tcolorbox}

\begin{attention}[Projet de fin d'année — Niveau professionnel attendu]
  Ce projet constitue le \textbf{projet de synthèse} de la formation BTS FED GCF.
  Le dossier technique doit atteindre un niveau proche d'une note de calcul professionnelle
  (style bureau d'études), avec plan de principe, notes de calcul justifiées, sélections
  de matériels sur catalogues réels et vérification réglementaire complète.
\end{attention}

\bigskip

% ================================================================
\section{Contexte professionnel}
% ================================================================

Vous êtes ingénieur d'études chez \textbf{FluidEnergy Bureau d'Études} (Bordeaux, 33).
Un promoteur, \textbf{Gironde Habitat}, vous confie la conception complète des installations
CVC, VMC et électriques CVC d'un \textbf{immeuble résidentiel R+4} neuf,
situé à \textbf{Bordeaux Bastide} (zone H2a, latitude 44,8°N).

Le programme impose des contraintes fortes sur le choix des équipements et le respect
de la RE2020 dans sa version résidentielle.

\subsection*{Description architecturale de l'immeuble}

\begin{figure}[htbp]
  \centering
  \begin{tikzpicture}[scale=0.9, >=Stealth, thick]

    % ============ COUPE TRANSVERSALE ============
    \node[font=\bfseries\small, FedBlue] at (6, 12.8) {Coupe transversale — Immeuble R+4 Gironde Habitat};

    % === Sous-sol parking ===
    \fill[gray!30] (0, -2.5) rectangle (12, 0);
    \draw[gray!60, thick] (0,-2.5) rectangle (12,0);
    \node[font=\small\bfseries, gray!70] at (6,-1.25) {Sous-sol — Parking 15 places};
    % PAC géo dans SS
    \draw[FedGreen, thick, fill=FedGreen!20, rounded corners=2pt]
      (0.5,-2.2) rectangle (3.0,-0.5);
    \node[FedGreen, font=\tiny\bfseries] at (1.75,-1.0) {PAC géoth.};
    \node[FedGreen, font=\tiny] at (1.75,-1.4) {2×25 kW};
    \node[FedGreen, font=\tiny] at (1.75,-1.7) {cascadées};
    % Sonde géo
    \draw[FedGreen, thick, ->] (1.75,-2.2) -- (1.75,-3.2)
      node[below, font=\tiny] {Sondes \SI{100}{\metre}};
    % Local technique SS
    \draw[FedBlue!60, thick, fill=FedBlue!10, rounded corners=2pt]
      (9.0,-2.2) rectangle (11.5,-0.5);
    \node[FedBlue, font=\tiny\bfseries] at (10.25,-1.0) {Local tech.};
    \node[FedBlue, font=\tiny] at (10.25,-1.4) {Chaufferie};
    \node[FedBlue, font=\tiny] at (10.25,-1.7) {+ Nourrice};

    % === RDC ===
    \fill[FedBlue!8] (0,0) rectangle (12,2.5);
    \draw[FedBlue!60] (0,0) rectangle (12,2.5);
    % Entrée + local vélos + poussettes
    \draw[brown!60, fill=brown!15] (4.5,0) rectangle (7.5,2.5);
    \node[brown!80, font=\tiny] at (6,1.25) {Hall + locaux communs};
    % 2 logements RDC
    \draw[FedBlue!40, fill=white] (0.2,0.2) rectangle (4.3,2.3);
    \node[FedBlue, font=\tiny] at (2.25,1.25) {T3 — \SI{62}{\metre\squared}};
    \draw[FedBlue!40, fill=white] (7.7,0.2) rectangle (11.8,2.3);
    \node[FedBlue, font=\tiny] at (9.75,1.25) {T3 — \SI{62}{\metre\squared}};
    \node[right, font=\tiny\bfseries, FedBlue] at (12.1,1.25) {RDC};

    % === R+1 à R+4 ===
    \foreach \niveau/\label in {2.5/R+1, 5.0/R+2, 7.5/R+3, 10.0/R+4}{
      \fill[FedBlue!5] (0,\niveau) rectangle (12,\niveau+2.3);
      \draw[FedBlue!40] (0,\niveau) rectangle (12,\niveau+2.3);
      % 4 logements par niveau
      \draw[FedBlue!30, fill=white] (0.2,\niveau+0.2) rectangle (2.8,\niveau+2.1);
      \draw[FedBlue!30, fill=white] (3.0,\niveau+0.2) rectangle (5.6,\niveau+2.1);
      \draw[FedBlue!30, fill=white] (6.4,\niveau+0.2) rectangle (9.0,\niveau+2.1);
      \draw[FedBlue!30, fill=white] (9.2,\niveau+0.2) rectangle (11.8,\niveau+2.1);
      \node[FedBlue, font=\tiny] at (1.5,\niveau+1.15) {T2};
      \node[FedBlue, font=\tiny] at (4.3,\niveau+1.15) {T3};
      \node[FedBlue, font=\tiny] at (7.7,\niveau+1.15) {T3};
      \node[FedBlue, font=\tiny] at (10.5,\niveau+1.15) {T2};
      \node[right, font=\tiny\bfseries, FedBlue] at (12.1,\niveau+1.15) {\label};
    }

    % === Toiture ===
    \fill[gray!20] (0,12.3) rectangle (12,12.8);
    \draw[gray, thick] (0,12.3) -- (12,12.3);
    % CTA toiture
    \draw[FedOrange, thick, fill=FedOrange!15, rounded corners=2pt]
      (4.5,12.3) rectangle (7.5,12.8);
    \node[FedOrange, font=\tiny\bfseries] at (6,12.55) {CTA VMC collective};

    % === Colonnes montantes (tirets) ===
    \draw[FedOrange, thick, dashed] (3.2, -0.5) -- (3.2, 12.3);
    \node[FedOrange, font=\tiny, rotate=90] at (3.35, 6) {Colonne chauffage};
    \draw[FedBlue, thick, dashed] (8.8, -0.5) -- (8.8, 12.3);
    \node[FedBlue, font=\tiny, rotate=90] at (8.65, 6) {Colonne ECS};

    % === Légende ===
    \draw[FedGreen, thick, fill=FedGreen!20] (0.2,-4.2) rectangle (1.2,-3.7);
    \node[right, font=\tiny] at (1.3,-3.95) {PAC géothermique};
    \draw[FedOrange, thick, dashed] (4.0,-3.95) -- (5.0,-3.95);
    \node[right, font=\tiny] at (5.1,-3.95) {Colonne chauffage};
    \draw[FedBlue, thick, dashed] (8.0,-3.95) -- (9.0,-3.95);
    \node[right, font=\tiny] at (9.1,-3.95) {Colonne ECS};

  \end{tikzpicture}
  \caption{Coupe transversale de l'immeuble R+4 Gironde Habitat — Bordeaux (H2a) —
    20 logements (2×T2 + 2×T3 par niveau R+1 à R+4, 2×T3 au RDC) —
    PAC géothermique en sous-sol + colonnes montantes + VMC hygroB}
  \label{fig:sr3-coupe-immeuble}
\end{figure}

\subsection*{Données du programme immobilier}

\begin{center}
\begin{tabular}{lrl}
  \toprule
  \textbf{Caractéristique} & \textbf{Valeur} & \textbf{Remarque} \\
  \midrule
  Nombre de logements & 20 & 4 T2 + 16 T3 (voir coupe) \\
  Surface totale SHON & \SI{1380}{\metre\squared} & RE2020 résidentiel \\
  Zones climatiques & H2a & $T_{ext,base} = \SI{-5}{\celsius}$ \\
  Hauteur sous plafond & \SI{2.5}{\metre} & Logements / \SI{3.0}{\metre} SS \\
  Sous-sol parking & 15 places & Ventilation mécanique imposée \\
  \midrule
  \textbf{Contraintes maître d'ouvrage} & & \\
  \midrule
  Production chaleur & PAC géothermique sondes & 2 unités cascadées \\
  Chauffage & Plancher chauffant basse T° & Chaque logement \\
  ECS & Chauffe-eau thermodynamique & 1 par logement \\
  VMC & Hygroréglable de type B & NF EN 13141-8 \\
  Tableau électrique CVC & 1 armoire dédiée par cage & Bilan de puissance obligatoire \\
  \bottomrule
\end{tabular}
\end{center}

% ================================================================
\section{Schéma général de l'installation}
% ================================================================

\begin{figure}[htbp]
  \centering
  \begin{tikzpicture}[
    >=Stealth, thick,
    bloc/.style={draw, rounded corners=3pt, minimum width=2.5cm, minimum height=1.0cm,
                 align=center, font=\scriptsize},
    dep/.style={FedOrange, thick, ->},
    ret/.style={FedBlue, thick, ->},
    geo/.style={FedGreen, thick, ->},
  ]

    % === SOURCES (gauche) ===
    \node[bloc, fill=FedGreen!20, draw=FedGreen, minimum height=1.4cm] (pac) at (0,4)
      {PAC géo\\2×25 kW\\COP = 4,2};
    \node[bloc, fill=brown!15, draw=brown!60] (sondes) at (0,1.5)
      {Sondes\\géothermiques\\2×100 m};
    \draw[geo, ->] (sondes.north) -- node[left, font=\tiny] {$T_{sol} = 13°C$} (pac.south);

    % === CIRCUIT PRIMAIRE ===
    \node[bloc, fill=FedBlue!10, draw=FedBlue] (tampon) at (4,4)
      {Ballon\\tampon\\500 L};
    \draw[dep] (pac.east) -- node[above, font=\tiny] {45°C} (tampon.west);
    \draw[ret] (tampon.west|-tampon.south) -- node[below, font=\tiny] {35°C}
      (pac.east|-pac.south);

    % === NOURRICE / COLLECTEUR ===
    \node[bloc, fill=FedOrange!15, draw=FedOrange, minimum width=3.0cm] (nourrice) at (8,4)
      {Nourrice\\départ/retour};
    \draw[dep] (tampon.east) -- (nourrice.west);
    \draw[ret, bend right=20] (nourrice.south) to (tampon.east|-tampon.south);

    % === CIRCUITS SECONDAIRES ===
    % PCH logements
    \node[bloc, fill=gray!10] (pch) at (12, 6)
      {Colonnes PCH\\20 logements\\$\phi = 25$°C / 35°C};
    % ECS
    \node[bloc, fill=cyan!10, draw=cyan!60] (ecs) at (12, 4)
      {CET\\20 logements\\200 L / logement};
    % Ventilation parking
    \node[bloc, fill=yellow!15, draw=yellow!60] (parking) at (12, 2)
      {VMC parking\\SS — \SI{3000}{\metre\cubed\per\hour}};

    \draw[dep, ->] (nourrice.north east) |- (pch.west);
    \draw[dep, ->] (nourrice.east) -- (ecs.west);
    \draw[dep, ->] (nourrice.south east) |- (parking.west);

    % Retours
    \draw[ret, ->] (pch.south) -- ++(0,-0.5) -| (nourrice.north);
    \draw[ret, ->] (ecs.south) -- ++(0,-0.3) -| (nourrice.north|-nourrice.south);
    \draw[ret, ->] (parking.north) -- ++(0,0.3) -| (nourrice.south);

    % === VASE D'EXPANSION ===
    \node[draw, circle, minimum size=0.7cm, fill=yellow!20, font=\tiny] (vase) at (4,6.5)
      {VE};
    \draw[gray, dashed] (vase) -- (tampon.north);
    \node[font=\tiny, gray] at (4,7.0) {25 L};

    % === COMPTAGE ===
    \node[draw, diamond, fill=purple!10, font=\tiny, minimum size=0.8cm] (compt) at (8,6.5)
      {Comptage\\calorimétrique};
    \draw[purple, dashed] (compt) -- (nourrice.north);

    % === RÉGULATION ===
    \node[bloc, fill=purple!10, draw=purple!50, minimum width=2.0cm] (reg) at (0,7.5)
      {Régulateur\\GTB\\BACnet IP};
    \draw[purple, dashed, ->] (reg.south) -- (pac.north);
    \draw[purple, dashed, ->] (reg.east) -| (nourrice.north|-compt) -- (compt.west);

    % === LÉGENDE ===
    \draw[dep] (0,-1.0) -- (1.2,-1.0) node[right, font=\tiny, black] {Départ chaud};
    \draw[ret] (3.5,-1.0) -- (4.7,-1.0) node[right, font=\tiny, black] {Retour froid};
    \draw[geo] (7.0,-1.0) -- (8.2,-1.0) node[right, font=\tiny, black] {Circuit géothermique};
    \draw[purple, dashed] (10.5,-1.0) -- (11.7,-1.0) node[right, font=\tiny, black] {GTB / Régulation};

  \end{tikzpicture}
  \caption{Schéma général de l'installation CVC — Immeuble R+4 Bordeaux —
    PAC géothermique 2×25 kW + ballon tampon 500 L + nourrice + colonnes PCH + CET + VMC parking}
  \label{fig:sr3-schema-installation}
\end{figure}

% ================================================================
\section{Cahier des charges}
% ================================================================

\begin{methode}[Livrables attendus — Note 1 : Dossier de conception complet]
\textbf{À remettre en semaine 8 — Minimum 50 pages — Niveau note de calcul professionnelle}

\begin{enumerate}
  \item \textbf{Analyse dimensionnelle et cohérence des calculs} (ch. 00) :
    \begin{itemize}
      \item Vérification dimensionnelle des 3 formules principales utilisées dans le dossier
      \item Tableau de synthèse des grandeurs, unités SI et ordres de grandeur attendus
    \end{itemize}

  \item \textbf{Bilan thermique complet} (ch. 01) :
    \begin{itemize}
      \item Calcul par type de logement (T2, T3) selon NF EN 12831
      \item $U_{bat}$ et comparaison seuil RE2020 résidentiel
      \item Déperditions totales bâtiment [\si{\kilo\watt}]
    \end{itemize}

  \item \textbf{PAC géothermique} (ch. 04) :
    \begin{itemize}
      \item Dimensionnement puissance, sélection modèle, calcul COP et SPF annuel
      \item Dimensionnement des sondes géothermiques (longueur, résistance thermique du sol)
      \item Vérification compatibilité avec le plancher chauffant (T° départ $\leq 45°C$)
    \end{itemize}

  \item \textbf{Réseau hydraulique} (ch. 02) :
    \begin{itemize}
      \item Dimensionnement des colonnes montantes et antennes de logement
      \item Équilibrage hydraulique par vannes d'équilibrage (méthode index + $k_{v}$)
      \item Vase d'expansion DTU 65.11 + sécurités (soupape, pressostat)
      \item Sélection des circulateurs (ErP 2015 + adaptation vitesse)
    \end{itemize}

  \item \textbf{Plancher chauffant} (ch. 05) :
    \begin{itemize}
      \item Calcul de la densité de flux $\phi$ [\si{\watt\per\metre\squared}] et du pas de pose
      \item Dimensionnement des boucles (longueur, $\Delta P$) par logement type
      \item Vérification de la température de surface $\leq 29°C$ (NF EN 1264)
    \end{itemize}

  \item \textbf{VMC hygroréglable de type B} (ch. 03) :
    \begin{itemize}
      \item Débits hygroréglables par local (cuisine, salle de bain, WC)
      \item Dimensionnement réseau collectif + gaine verticale en sous-pression
      \item SFP réseau [\si{\watt\per\litre\per\second}] et classe énergétique
    \end{itemize}

  \item \textbf{Étude acoustique CVC} (ch. 09) :
    \begin{itemize}
      \item Niveau de puissance acoustique $L_w$ des équipements (PAC, VMC, circulateurs)
      \item Vérification du niveau NC dans les logements (limite NRA : $\leq 30 dB(A)$ nuit)
      \item Mesures constructives si dépassement (silencieux, plots antivibratoires)
    \end{itemize}

  \item \textbf{Électrotechnique CVC} (ch. 10) :
    \begin{itemize}
      \item Bilan de puissance installée et absorbée [\si{\kilo\watt}]
      \item Schéma unifilaire de l'armoire CVC (disjoncteurs, contacteurs, variateurs)
      \item Calcul de $\cos\varphi$ global et compensation éventuelle
    \end{itemize}

  \item \textbf{Bilan RE2020 complet} (ch. 08) :
    \begin{itemize}
      \item Calcul Bbio, Cep\textsubscript{nr}, Ic\textsubscript{énergie}
      \item Attestation de prise en compte RE2020 (simulation Pleiades ou justification)
      \item Vérification des 3 seuils réglementaires
    \end{itemize}

  \item \textbf{Budget et planning} :
    \begin{itemize}
      \item Devis estimatif par lot (génie thermique, VMC, électricité CVC)
      \item Planning de chantier Gantt sur 8 semaines (gros œuvre terminé)
    \end{itemize}
\end{enumerate}
\end{methode}

\begin{methode}[Livrables attendus — Note 2 : Soutenance devant commission]
\textbf{Durée : 30 min exposé + 15 min questions — Commission mixte : enseignant + professionnel invité}

\begin{itemize}
  \item Plan de principe A3 imprimé à apporter (schéma général de l'installation)
  \item Présentation des choix technologiques et comparaison avec 2 alternatives justifiées :
    \begin{itemize}
      \item Alternative A : PAC géo vs PAC air/eau (COP, investissement, faisabilité)
      \item Alternative B : PCH vs radiateurs basse T° (inertie, confort, coût)
    \end{itemize}
  \item Questions de jury sur \textbf{n'importe quel calcul du dossier} — préparation individuelle obligatoire
  \item Mise en situation : le jury annonce une contrainte de chantier (ex : sondes géo impossibles) — adapter le projet en 5 minutes
\end{itemize}
\end{methode}

% ================================================================
\section{Grille d'évaluation}
% ================================================================

\begin{center}
\begin{tabular}{p{6cm}ccp{4cm}}
  \toprule
  \textbf{Critère} & \textbf{Barème} & \textbf{Note} & \textbf{Indicateurs} \\
  \midrule
  \multicolumn{4}{l}{\textbf{Note 1 — Dossier de conception (20 pts)}} \\
  \midrule
  Bilan thermique + PAC géo & 4 pts & /4 & Puissances cohérentes, SPF calculé \\
  Réseau hydraulique + PCH & 4 pts & /4 & Équilibrage, $\Delta P$, vase \\
  VMC hygroB + acoustique & 3 pts & /3 & SFP, NC logements vérifiés \\
  Électrotechnique CVC & 2 pts & /2 & Armoire, $\cos\varphi$ \\
  RE2020 complet & 3 pts & /3 & 3 indicateurs, seuils respectés \\
  Budget + planning & 2 pts & /2 & Cohérence €, Gantt structuré \\
  Qualité professionnelle & 2 pts & /2 & Mise en page, schémas, vocabulaire \\
  \midrule
  \multicolumn{4}{l}{\textbf{Note 2 — Soutenance commission (20 pts)}} \\
  \midrule
  Exposé (clarté, structure) & 5 pts & /5 & Plan logique, gestion du temps \\
  Maîtrise des calculs (jury) & 7 pts & /7 & Réponses exactes aux questions \\
  Argumentation des choix & 5 pts & /5 & Alternatives comparées, pertinence \\
  Réactivité situation imposée & 3 pts & /3 & Adaptation rapide et cohérente \\
  \bottomrule
\end{tabular}
\end{center}

% ================================================================
\section{Ressources documentaires}
% ================================================================

\begin{itemize}
  \item \textbf{Normes :} NF EN 12831, NF EN 14511 (PAC), NF EN 1264 (PCH), NF EN 13141-8 (VMC hygro B)
  \item \textbf{DTU :} 65.11 (sécurités hydrauliques), 68.1 (VMC), 60.1 (plomberie)
  \item \textbf{Réglementation :} RE2020 résidentiel (arrêté 04/08/2021), Décret F-Gas 2024
  \item \textbf{Logiciels conseillés :} Pleiades+Comfie (RE2020), ClimaWin, Trace 700
  \item \textbf{Catalogues :} Kensa (PAC géo), Grundfos (circulateurs), Zehnder (VMC), Giacomini (PCH)
  \item \textbf{Données météo Bordeaux :} Fichier DRY METEONORM — station 33522 — 2 900 DJU
  \item \textbf{Cours :} Tous les chapitres 00 à 11 du référentiel GCF
\end{itemize}

\finchapitre
